\chapter{Assembly Languages, Expressions, Symbols, and Sections}
\label{ch:assembly-language}

A compiler can decide that a function needs an integer addition without knowing the address at which that addition will execute.
The assembler preserves this separation.
It translates a target-specific notation into instruction and data bytes, but it also records unresolved relationships between those bytes and named locations.
If it replaces an unknown address by an unexplained zero, the linker cannot recover the missing relationship.
If it insists that every address be final immediately, separately compiled objects become impossible.
The useful output is therefore both a byte sequence and a structured account of what remains undecided.

The next four chapters develop that output from the source language upward.
The working target is a deliberately documented subset of x86-64, with ELF64 relocatable objects and Lua preprocessing.
The principles apply to other instruction sets, but instruction sizes, relocation arithmetic, and ABI conventions always belong to a particular target.
Salomon's treatment of assembler organization supplies useful background; the ELF specification supplies the file-level contract~\cite{salomon-1993-assemblers,xinuos-2026-elf}.

\section{The assembler has several clients}
\label{sec:ch53-clients}

A programmer wants readable mnemonics, symbolic names, diagnostics, and predictable expansions.
A compiler wants stable directives for sections, alignment, linkage, debugging, and unwind information.
A linker wants machine-readable definitions, references, and fixups.
The processor wants only correctly encoded instructions at execution time.
These clients do not require the same representation.
An assembler listing may show a symbolic call even though the object contains a placeholder field and a relocation record rather than a final call displacement.

Figure~\ref{fig:ch53-contracts} distinguishes the source, object, and execution boundaries.
The source-to-object arrow is not a commitment to final virtual addresses.
Conversely, a linker does not normally parse the source mnemonic to discover what a relocation means; the object must already carry that meaning.
This is why an assembler cannot be implemented as a dictionary replacing each mnemonic with one fixed byte string.
Operands, widths, instruction forms, and location-dependent fields determine the encoding.

\begin{figure}[htbp]
\centering
\begin{tikzpicture}[node distance=15mm]
\node[basic block] (src) {mnemonics + directives};
\node[basic block,below=of src] (obj) {bytes + symbols + relocations};
\node[basic block,below=of obj] (image) {mapped executable bytes};
\draw[cfg edge] (src)--node[right,font=\small] {assembler}(obj);
\draw[cfg edge] (obj)--node[right,font=\small] {linker and loader}(image);
\node[font=\small,left=10mm of obj,align=right] {addresses may\\remain unknown};
\end{tikzpicture}
\caption{The object boundary preserves unresolved address relationships instead of pretending every byte is final.}
\label{fig:ch53-contracts}
\end{figure}

A \term{directive} changes assembler state or emits structured data rather than directly naming a processor instruction.
Selecting a section, declaring a global symbol, and requesting alignment are directives.
Some directives emit no bytes but still change the object substantially.
Marking a symbol global can make the difference between a private implementation label and a definition that satisfies another object's reference.

\section{Choose a source grammar before choosing an encoder}
\label{sec:ch53-syntax}

Assembly languages are not interchangeable just because they name the same machine.
Operand order, register prefixes, immediate syntax, memory addressing, comments, and directive names differ between dialects.
The case study uses destination-first operands, bare register names, semicolon comments, and one statement per generated line.
It accepts an optional label before a statement.
Its \code{lea rsi, message} is a restricted notation for RIP-relative address formation, not a promise to parse every Intel-syntax memory operand.
The restricted grammar is part of the public interface.

Tokenization must not let a malformed operand quietly become an unknown symbol.
For example, a register-width mismatch in \code{mov rax, eax} is an encoding error in this interface, not an external reference named \code{eax}.
An unsupported register such as \code{r8} is rejected rather than partially encoded using the low three bits of its register number.
Extended-register support requires the corresponding REX bits in every relevant operand position.

An implementation can parse into statement records containing an opcode identifier, typed operands, a source span, and a chosen encoding form.
That representation makes later layout passes independent of textual rescanning.
The small case study combines parsing and fixed-size emission, retaining symbolic fixups separately.
Chapter~\ref{ch:assembly-preprocessing} explains when variable-size encoding makes a richer fragment representation necessary.

\section{A location counter belongs to a section}
\label{sec:ch53-counters}

A \term{location counter} records the next offset in one output section.
Switching from text to data does not discard the text counter.
Returning to text continues from its previous offset unless the language explicitly specifies another behavior.
Table~\ref{tab:ch53-counters} works a sequence containing two instruction sizes, a data word, and alignment.
All numbers are section-relative; no final address has been assigned.

\begin{table}[htbp]
\caption{Section switches preserve independent counters; alignment contributes real padding.}
\label{tab:ch53-counters}
\centering
\begin{tabular}{llll}
\toprule
Operation & Text counter & Data counter & New definition \\
\midrule
Start & $0$ & $0$ & none \\
Define start; emit 5 bytes & $5$ & $0$ & start: text $+0$ \\
Switch data; emit 8 bytes & $5$ & $8$ & cell: data $+0$ \\
Switch text; align to 8 & $8$ & $8$ & three padding bytes \\
Define next; emit 2 bytes & $10$ & $8$ & next: text $+8$ \\
\bottomrule
\end{tabular}
\end{table}

Figure~\ref{fig:ch53-counters} draws the holes and contributions separately.
Padding inside a text section is still part of that section's size.
If control can fall through the padding, its byte values must have a defined executable meaning.
Our assembler's alignment directive emits zeros in file-backed sections, so code must jump around such padding or use it only between independently entered code blocks.
Production assemblers often choose target-specific no-operation sequences for code alignment; that is an encoding policy rather than a mathematical property of alignment.

\begin{figure}[htbp]
\centering
\begin{tikzpicture}[x=9mm,y=11mm]
\draw (0,0) rectangle (5,1); \node at (2.5,0.5) {5 bytes};
\draw[dashed] (5,0) rectangle (8,1); \node at (6.5,0.5) {pad};
\draw (8,0) rectangle (10,1); \node at (9,0.5) {2};
\node[anchor=east] at (-0.2,0.5) {text};
\draw (0,-2) rectangle (8,-1); \node at (4,-1.5) {8-byte cell};
\node[anchor=east] at (-0.2,-1.5) {data};
\foreach \x in {0,5,8,10} \node[font=\small] at (\x,1.4) {\x};
\end{tikzpicture}
\caption{Text and data offsets live in separate coordinate systems even when their numeric values coincide.}
\label{fig:ch53-counters}
\end{figure}

\section{Represent expressions by what can be relocated}
\label{sec:ch53-expressions}

An integer expression and an address expression have different relocation behavior.
The absolute value $24$ stays $24$ when a section moves.
The address of a label at text offset $24$ gains the text base.
Treating both as an untagged integer loses information needed by the linker.
A useful internal value domain contains absolute integers, section-relative values, unresolved symbol-plus-addend expressions, and invalid expressions.

For section bases $B_s$, represent a section-relative value as $B_s+c$.
Adding an absolute $k$ gives $B_s+(c+k)$.
Subtracting two values from the same section cancels the common base and gives an absolute difference.
Adding two section-relative addresses generally produces $B_s+B_t+c$, which a simple one-symbol relocation cannot represent.
Different object formats support different expression languages; acceptance must follow the chosen format rather than the host evaluator's ability to compute arbitrary arithmetic.
Table~\ref{tab:ch53-expressions} lists representative cases.

\begin{table}[htbp]
\caption{Expression classification preserves relocation dependence instead of erasing it during evaluation.}
\label{tab:ch53-expressions}
\centering
\begin{tabular}{lll}
\toprule
Expression & Classification & Condition \\
\midrule
$17+6$ & absolute $23$ & checked integer arithmetic \\
$(B_t+12)+4$ & text-relative $B_t+16$ & same base retained \\
$(B_t+28)-(B_t+12)$ & absolute $16$ & both offsets final \\
$(B_t+12)-(B_d+4)$ & unresolved base difference & needs richer contract \\
$external+8$ & symbol plus addend & supported relocation kind \\
\bottomrule
\end{tabular}
\end{table}

\begin{algorithm}[htbp]
\caption{Classify addition and subtraction in a simple relocatable-expression domain.}
\label{alg:ch53-expressions}
\begin{pseudocode}
\Require tagged operands with checked offsets; operator plus or minus
\Ensure absolute result, one-base result, or explicit unsupported-expression error
\If{both operands are absolute}
  \State compute the checked integer operation
\ElsIf{one-base left operand and absolute right operand}
  \State preserve the base and adjust its offset with checked arithmetic
\ElsIf{absolute left plus one-base right}
  \State preserve the right base and add the absolute offset
\ElsIf{subtracting two final offsets with the same base}
  \State return their checked absolute difference
\Else
  \State reject or retain a richer symbolic expression supported by the object format
\EndIf
\end{pseudocode}
\end{algorithm}

Algorithm~\ref{alg:ch53-expressions} has constant work per operator once symbols are classified, but determining whether an offset is final may depend on layout convergence.
If instruction relaxation can change the distance between two labels, their same-section difference is not ready to freeze merely because the section names match.
The case study avoids that complication by using fixed encoding sizes and accepting only an integer or one symbol plus an integer addend in emitted fields.
It does not implement general assembler expression algebra.

\section{Symbols have identities and states}
\label{sec:ch53-symbols}

A symbol table entry needs more than a spelling and a number.
Useful attributes include definition state, binding, section identity, section offset, optional size/type, and the source location of the defining declaration.
An unresolved reference creates an entry whose value is not yet known.
A later label definition updates that entry rather than creating another symbol with the same spelling.
Two definitions of the same local name in one object are an error; the same local spelling in two different objects is normally harmless.

Figure~\ref{fig:ch53-symbol-states} separates declaration from definition.
The assembler can require explicit external declarations, or it can export every unresolved name as undefined; both policies are possible.
The case study chooses explicit \code{.extern} or \code{.weak} declarations so that misspelled private labels do not silently turn into link-time dependencies.
This is stricter than many production assemblers and is documented as such.

\begin{figure}[htbp]
\centering
\begin{tikzpicture}[node distance=17mm and 24mm]
\node[basic block] (u) {unseen name};
\node[basic block,right=of u] (r) {referenced, undefined};
\node[basic block,below=of r] (d) {defined: section + offset};
\node[basic block,left=of d] (e) {external declaration};
\draw[cfg edge] (u)--(r); \draw[cfg edge] (r)--node[right,font=\small] {label}(d);
\draw[cfg edge] (r)--(e); \draw[cfg edge] (e)--node[below=3mm,font=\small] {optional definition}(d);
\node[basic block,below=of d] (err) {duplicate-definition error};
\draw[back edge] (d)--node[right,font=\small] {second label}(err);
\end{tikzpicture}
\caption{An external declaration changes linkage without inventing a definition or assigning address zero.}
\label{fig:ch53-symbol-states}
\end{figure}

Stable symbol identities matter when an object writer reorders symbols.
ELF requires local symbols before nonlocals; source order need not satisfy that condition~\cite{xinuos-2026-elf}.
Internal fixups should therefore reference an internal symbol identifier.
Only final serialization translates that identifier into an ELF symbol-table index.
Reordering the symbol table without translating relocation references yields well-formed records that name the wrong entities.

\section{Forward references do not dictate the number of source passes}
\label{sec:ch53-passes}

A classic two-pass assembler first determines sizes and labels, then emits encoded bytes using the completed symbol table.
That organization is useful when all sizes can be determined from syntax before values are known.
It is not a theorem that a one-pass source reader cannot emit relocatable objects.
A source reader can emit known bytes, retain fixups in memory, and serialize unresolved relationships after reaching the end of the source.
The number of textual passes, number of internal layout iterations, and kind of output object are separate choices.
Historical descriptions of particular one-pass loaders should not be generalized into restrictions on all assemblers~\cite{salomon-1993-assemblers}.

\begin{algorithm}[htbp]
\caption{Read fixed-size assembly once while retaining symbolic fixups.}
\label{alg:ch53-fixups}
\begin{pseudocode}
\Require parsed statements whose instruction sizes do not depend on symbol values
\Ensure section bytes, definitions, and unresolved field relationships
\ForAll{statements in source order}
  \If{the statement defines a label}
    \State bind its stable symbol identifier to current section and offset
  \EndIf
  \If{the statement emits an instruction or data field}
    \State append known bytes and reserve the exact size of each unresolved field
    \State record field offset, symbol identifier, addend, and relocation kind
  \Else
    \State apply the directive to section or symbol state
  \EndIf
\EndFor
\State diagnose unresolved private names and conflicting definitions
\State serialize bytes, symbols, and translated relocation references
\end{pseudocode}
\end{algorithm}

With $N$ statements, $B$ emitted bytes, and $R$ fixups, this organization needs $O(N+B+R)$ work plus symbol lookup and expression evaluation.
A linear symbol array makes repeated lookup quadratic in the number of distinct names in the worst case.
Hashing gives expected constant-time lookup under its distribution assumptions; balanced trees give logarithmic worst-case lookup.
The case study intentionally uses bounded linear tables for reviewability, and does not claim production-scale linear performance.

\section{Storage directives are not all file bytes}
\label{sec:ch53-storage}

Initialized data has both a file representation and a memory extent.
Zero-initialized storage can have only a memory extent.
In ELF, \code{SHT\_NOBITS} records size and alignment without storing that many payload bytes~\cite{xinuos-2026-elf}.
A label inside such a section still has a meaningful section-relative offset.
Its address can be relocated into a pointer elsewhere even though the labeled storage has no initial file payload.

Table~\ref{tab:ch53-storage} distinguishes three quantities that are easy to conflate.
A large zero-fill request must increase memory size and later labels without forcing the assembler to allocate an equally large byte array.
The implementation nevertheless caps the requested extent, because a compact object can otherwise describe an arbitrarily large eventual allocation.

\begin{table}[htbp]
\caption{Zero-filled storage contributes to layout without contributing equal file payload.}
\label{tab:ch53-storage}
\centering
\begin{tabular}{llll}
\toprule
Contribution & File bytes & Memory bytes & Label step \\
\midrule
64 initialized bytes & $64$ & $64$ & $64$ \\
64 bytes in NOBITS & $0$ & $64$ & $64$ \\
Alignment from 17 to 32 & depends on section type & $15$ & $15$ \\
Binding declaration & $0$ payload & $0$ & $0$ \\
\bottomrule
\end{tabular}
\end{table}

\section{Diagnostics should preserve the failed decision}
\label{sec:ch53-diagnostics}

Report whether a failure came from tokenization, unsupported syntax, expression classification, instruction encoding, or unresolved linkage.
An out-of-range immediate is different from an unsupported addressing mode even if both occur in the second operand.
Include the operand's spelling and source location before lowering destroys that context.
For generated assembly, distinguish the generated line number from the Lua line that produced it.
The case study reports the former for assembly errors and Lua's own diagnostic for preprocessing errors; a production macro system should retain an expansion stack as well.

Do not write the destination object until semantic checks and serialization have succeeded.
This keeps an invalid source from replacing a previously useful object with a partial header.
An I/O failure during the final write remains a separate issue; atomic replacement through a temporary file is a useful production extension.
The small tools remove a partially written output on a detected write failure.

\section{Exercises}
\label{sec:ch53-exercises}
\begin{enumerate}
\item\label{ex:ch53-counter} Extend Table~\ref{tab:ch53-counters} by switching to data, aligning to 16, defining tail, and emitting four bytes. Give tail's value and both final counters.
\item\label{ex:ch53-expression} Classify $(B_t+40)-(B_t+9)$ and $(B_t+40)-(B_d+9)$. State the extra assumption needed before freezing the first result.
\item\label{ex:ch53-symbol-index} Internal symbols are created in order global g, local a, local b. With ELF null entry zero, give the serialized indices and the mapping required by a fixup to g.
\item\label{ex:ch53-nobits} A NOBITS section contains 24 bytes, then alignment to 32, then 12 bytes. Compute file payload, section extent, and the last label's offset.
\item\label{ex:ch53-forward} Explain why a fixed five-byte call to a later label does not force a second source-text pass, and identify the record that must survive until object emission.
\end{enumerate}

\section{Further reading}
\label{sec:ch53-reading}
Salomon develops directives, symbols, macros, and assembler organizations across several historical machines~\cite{salomon-1993-assemblers}.
The ELF specification is the normative reference for section types, symbol binding, and symbol-table ordering~\cite{xinuos-2026-elf}.
The AMD64 ABI supplement connects those generic structures to x86-64 address conventions~\cite{matz-2012-amd64}.
